\section{World Money, Financial Subordination, and the Peripheral Central Bank Balance Sheet}
\label{sec:macro_framework}

This section develops the theoretical framework linking the global currency hierarchy to domestic macroeconomic destabilization. The analysis traces a clear analytical sequence: the asymmetric structure of world money generates peripheral settlement constraints, which concentrate balance-of-payments vulnerability directly onto the central bank balance sheet. In response to structural bottlenecks and foreign-exchange scarcity, domestic monetary issuance functions as an endogenous accommodation mechanism rather than an autonomous policy choice. The resulting balance-sheet strain alters macroeconomic transmission, generating state-dependent responses that are subsequently estimated through dynamic threshold modeling.

\subsection{World Money, the Universal Equivalent, and the Monetary Expression of Value}
\label{sec:B1_worldmoney_mev}

\textbf{3.1} In Marxian political economy, money originates as the commodity form of the universal equivalent: the unique social medium whose function is to express the value of all commodities and serve as the general means of payment \citep[\S3.1]{Basu2022}. Gold historically occupied this role not through inherent natural properties, but through a social process of market validation, acting as the bearer of abstract human labor. 

\textbf{3.2} The historical transit to credit money transforms this institutional form without dissolving its underlying function. As banking systems developed to intermediate commercial credit, central bank liabilities emerged as the legal anchor of domestic circulation \citep[\S3.2.3]{Basu2022}. The commodity anchor disappeared from domestic circulation, but the imperative of social validation remained. Central bank liabilities and bank deposits perform the role that bullion once fulfilled: providing the monetary expression through which private, decentralized labor is validated as social labor. Following the Monetary Expression of Value (MEV) framework \citep{Foley1982,Basu2022}, the value of money links aggregate labor-time to the circulating monetary unit within a national formation.

\textbf{3.3} Once circulation crosses national borders, however, credit money confronts the world market. While Marx presented world money in the physical form of bullion shedding national garbs \citep[p.~51]{Ivanova2013}, a global credit-money system reproduces the universal equivalent through an international currency hierarchy \citep{Vasudevan2009,Ivanova2020}. National currencies are stratified by their degree of international acceptability, liquidity, and convertibility. World money is the historically specific form wherein the liabilities of the dominant hegemonic power operate as the universal means of international payment, reserve asset, and unit of account.

\subsection{The Bretton Woods System and Peripheral Balance-of-Payments Constraints}
\label{sec:B2_BrettonWoods_BOP}

\textbf{3.4} The post-war Bretton Woods monetary system codified this currency hierarchy institutionally. By pegging the U.S. dollar to gold at \$35 per ounce while anchoring other currencies to the dollar, the international monetary order formalized the dollar as the operative form of world money \citep[p.~483]{Vasudevan2009}. This arrangement granted the hegemonic center an \textit{``exorbitant privilege''}: the ability to borrow internationally in its own currency, finance external deficits without reserve constraints, and act as banker to the world \citep[pp.~243--244]{Vasudevan2019}.

\textbf{3.5} For peripheral economies, this institutional structure generated an acute structural asymmetry. Peripheral states cannot settle international accounts in their own domestic currencies; external liabilities and essential imports must be discharged in world money. The peripheral circuit of capital confronts an absolute balance-of-payments constraint \citep{Vasudevan2016,Abeles2013}. Domestic credit expansion is bounded by the capacity to command foreign exchange through export earnings, terms of trade, and access to international credit lines. Reserve accumulation constitutes an externally enforced structural prerequisite of economic reproduction \citep[p.~244]{Vasudevan2019}.

\textbf{3.6} The unilateral closure of the gold window by the United States on August 15, 1971 shattered the fixed-parity anchor of Bretton Woods agreements leading to their demise. While this rupture removed gold convertibility, it preserved dollar centrality, ushering in an era of global floating fiat money \citep[p.~486]{Vasudevan2009}. For the nascent socialist project of the Unidad Popular, this systemic shock dismantled the international monetary stability upon which its initial macroeconomic projections rested.

\subsection{The Peripheral Central Bank Balance Sheet: From Reserve Adequacy to Balance-Sheet Solvency}
\label{sec:B3_levr}

\textbf{3.7} Financial subordination materializes concretely on the balance sheet of the peripheral central bank. Mainstream open-economy models interpret this balance sheet through the lens of ``\textit{reserve adequacy}'' and precautionary self-insurance, treating international reserves as an inventory buffer optimized against friction costs to insulate the economy from sudden stops \citep{Jeanne2007,CespedesChang2024}. This normative framing assumes that reserve accumulation represents the voluntary, unconstrained choice of a sovereign optimizing agent.

\textbf{3.8} In peripheral capitalism, reserve trajectories are the structural outcome of the balance-of-payments constraint, unequal exchange, and financial subordination \citep{Vasudevan2009,Abeles2013}. The two sides of the peripheral central bank's balance sheet reflect fundamentally distinct monetary forms. On the liability side, domestic money represents credit obligations denominated in the national unit of account, which the central bank can issue to meet payments-system liquidity \citep[\S3.2.3]{Basu2022}. On the asset side, gross international reserves ($IR_t$, measured in USD millions) consist of liquid external claims denominated in world money---convertible foreign exchange balances, short-term foreign sovereign assets, and monetary gold---which the peripheral state cannot emit and can acquire only through export surpluses, net foreign capital inflows, or international credit lines. While the central bank can expand domestic liquidity endogenously to accommodate internal clearing, its ability to settle external transactions and service foreign debt depends strictly on this stock of world money. 

\textbf{3.9} This balance-sheet structure is governed by the Central Bank Solvency Ratio:
\begin{equation}
	\text{SolvR}_t \;\equiv\; \frac{e_t \cdot IR_t}{M_t} \;\equiv\; \frac{F_t}{M_t}
	\label{eq:solvency_ratio}
\end{equation}
where $IR_t$ denotes international reserves in USD, $e_t$ is the nominal exchange rate (Escudos/Pesos per USD), $F_t \equiv e_t \cdot IR_t$ is the domestic-currency value of foreign reserve assets, and $M_t$ represents the generic theoretical concept of domestic monetary liabilities.\footnote{The Solvency Ratio is the inverse of the Leverage Ratio ($\text{LevR}_t \equiv M_t / (e_t \cdot IR_t) = \text{SolvR}_t^{-1}$). The solvency formulation is preferred throughout this study because as reserves shrink toward zero, the leverage ratio explodes and distorts the statistical distribution, whereas the Solvency Ratio converges smoothly to zero, yielding a well-behaved index for threshold estimation.} In the empirical baseline developed in Section~\ref{sec:data_sources} and Section~\ref{sec:historical_empirical_results}, central bank liabilities are operationalized with high-powered base money $H_t$ ($\text{SolvR}^H_t \equiv [e_t \cdot IR_t] / [H_t \cdot 1000]$), while narrow commercial-bank money ($M1_t$) is reserved for the banking-sector deconstruction.

In foreign-currency units, the reserve-flow identity is $\Delta IR_t = CA_t + KA_t + EO_t + SDR_t$, where $CA_t$, $KA_t$, $EO_t$, and $SDR_t$ denote current account receipts, capital account net transactions, errors and omissions, and SDR allocations. Because $F_t \equiv e_t \cdot IR_t$, the exact discrete change in domestic-currency foreign assets is $\Delta F_t = e_t \Delta IR_t + IR_{t-1} \Delta e_t$, distinguishing reserve flows from exchange-rate valuation effects. Following the formal balance-sheet accounting derived in Appendix~\ref{app:bop_levr}, the exact discrete law of motion governing the level Solvency Ratio is:
\begin{equation}
	\Delta \text{SolvR}_t \;=\; \frac{\Delta F_t \;-\; \text{SolvR}_{t-1} \Delta M_t}{M_t}
	\label{eq:solvency_law_motion}
\end{equation}
where the numerator decomposes the net change in domestic-currency foreign assets ($\Delta F_t$) and the domestic liquidity dilution effect evaluated at the lagged solvency stock ($\text{SolvR}_{t-1} \Delta M_t$), scaled by current monetary liabilities $M_t$.

Equation~\ref{eq:solvency_law_motion} highlights the analytical distinction between the level stock of foreign-exchange backing ($\text{SolvR}_t$) and the dynamic flow gap. While $\text{SolvR}_t$ tracks cumulative external backing, the stationary transition variable estimated in the Threshold VAR (Section~\ref{sec:tvar_solvency_regimes}) is the predetermined Solvency Growth Gap:
\begin{equation}
	\Delta s_{t-1} \;\equiv\; g^F_{t-1} - g^H_{t-1} \;=\; \Delta \ln \text{SolvR}^H_{t-1} \times 100,
	\label{eq:solvency_growth_gap_def}
\end{equation}
with $g^F_t \equiv 100 \Delta \ln F_t$ and $g^H_t \equiv 100 \Delta \ln H_t$. The estimated threshold ($\hat{\gamma} = -4.615\%$) identifies a structural break in the growth rate of foreign assets relative to base money, rather than a threshold in the level stock of reserves. The complete accounting derivation is set out in Appendix~\ref{app:bop_levr}.

\subsection{Money Endogeneity vs. Exogeneity and the Peripheral Paradox}
\label{sec:B4_money_endogeneity_paradox}
\label{sec:money_endogeneity_paradox}

\textbf{3.10} Grounding the peripheral balance sheet requires resolving the foundational debate over the causality between bank credit, deposits, and central bank reserves \citep{RochonRossi2013,Lavoie2014}. Mainstream monetarism and orthodox macroeconomics operate under Fractional Reserve Theory (FRT) and loanable-funds: the central bank exogenously fixes the monetary base ($H_t$), which mechanically dictates bank lending and broad money through a stable fractional-reserve multiplier ($H \to \text{Deposits} \to \text{Loans} \to P$). In sharp contrast, Post-Keynesian Endogenous Money Theory (EMT) and Marxian monetary circuits establish that real-world banking operates in reverse causal sequence \citep{Kaldor1982,Lavoie2014}:
\begin{equation}
	\Delta \text{Bank Loans} \;\longrightarrow\; \Delta \text{Deposits} \;\longrightarrow\; \Delta \text{Central Bank Reserves} (H)
	\label{eq:endogenous_money_sequence}
\end{equation}
Commercial banks extend working-capital loans to firms to finance payrolls and production. These loans simultaneously create bank deposits. The central bank, far from controlling the money supply as an exogenous policy instrument, is compelled to accommodate required reserves ex-post at the policy discount rate to maintain clearance in the interbank payments system. As \citet{Kaldor1982} argued, credit money has no independent supply function: it comes into existence through the demand for bank lending and is extinguished through repayment. The central bank must accommodate this process to prevent the collapse of the credit and payments system. Monetarism, in this light, is not a neutral technical rule but a political strategy that deploys high interest rates and monetary contraction to induce unemployment and break trade union bargaining power. Under EMT, the traditional money multiplier ($mm_t \equiv M_t / H_t$) is not an exogenous causal driver but an unstable, ex-post accounting ratio reflecting commercial bank liquidity preferences, cash-drain ratios, and portfolio balance constraints. In structuralist endogenous-money approaches, credit supply is upward-sloping rather than horizontal \citep{Dow2006,Rochon1999}. In peripheral capitalism, however, commercial banks' capacity and willingness to intermediate credit are bounded by external liquidity constraints and foreign exchange exposure.

\textbf{3.11} This institutional reality establishes what I define as the \emph{Peripheral Paradox}. In core economies issuing world money, the central bank accommodates domestic credit without facing an absolute currency constraint. In the periphery, however, while domestic credit generation is naturally endogenous ($\text{Loans} \to M_t$), the ultimate means of international settlement is bound to an exogenously constrained reserve asset ($e_t \cdot IR_t$). The peripheral central bank must back an endogenously expanding domestic credit pyramid with an externally constrained stock of world money. When external credit embargoes, copper price collapses, or private capital flight drain foreign reserves, the state cannot act as an unconstrained lender of last resort in hard currency; it must choose between contracting domestic activity or suffering balance-sheet insolvency. Under the UP, the nationalized commercial banking system under the State Development Corporation (CORFO) extended defensive working-capital credit to socialized and private enterprises facing distribution bottlenecks and price controls, compelling the Central Bank to monetize enterprise deficits. Domestic money growth thus functioned as an endogenous accommodation valve, driven by structural cost-push pressures rather than discretionary printing press impulses. This reverse causal sequence ($\Delta \text{Loans} \to \Delta \text{Deposits} \to \Delta H$) and its balance-sheet constraint under the Peripheral Paradox provide the operational hypotheses evaluated in the sequential Granger causality battery in Section~\ref{sec:granger_tournament}.

\subsection{Peripheral Financial Fragility and Central-Bank Solvency}
\label{sec:B5_foley_chamorro_minsky}

\textbf{3.12} To understand how this balance-sheet contradiction generates macroeconomic crisis, I draw on the financial fragility framework of \citet{Minsky1986}. In Minsky's canonical framework, financial fragility is endogenously generated during prolonged expansions: stability breeds instability as private firms and banks increase leverage, increasing their vulnerability to liquidity shortfalls \citep{Kregel2008}. However, Minsky formulated this dynamic for private corporate actors in a closed, metropolitan economy possessing complete monetary sovereignty, where an unconstrained central bank (the ``Big Bank'') and fiscal authority (the ``Big Government'') stabilize an intrinsically unstable system. Minsky's categorical taxonomy of hedge, speculative, and Ponzi finance was developed specifically for private cash flows and corporate debt commitments. The peripheral central bank is neither private nor unconstrained: it does not lever up out of speculative euphoria, but is compelled to intermediate liquidity under structural balance-of-payments constraints. Rather than attempting to map Minsky's categorical taxonomy directly onto the central bank, this study retains Minsky's foundational intuition of financial fragility while reinterpreting its institutional locus through peripheral open-economy conditions.

\textbf{3.13} Reinterpreting financial fragility for the peripheral central bank rests on two complementary foundations. First, \citet[p.~21]{Foley2003} demonstrates that during acute financial distress, the state and the central bank as lender of last resort are compelled to absorb the insolvency of private enterprises onto their own balance sheets. The central bank expands credit not out of profit-seeking, but to prevent systemic payments breakdown and societal disruption. Second, \citet{Chamorro2025} demonstrates that financial instability in peripheral Latin American economies operates as an \emph{externally amplified endogenous fragility}. Minsky's dynamic cannot be applied to dependent formations without recognizing that external shocks from the center---the collapse of world monetary arrangements, terms-of-trade declines, and credit embargoes---serve as structural triggers that detonate domestic balance-sheet fragility through the international currency hierarchy. During acute external instability, the peripheral central bank enters a state of defensive accommodation under severe balance-of-payments constraints. When external restrictions cut off commercial trade credit, refusing domestic liquidity expansion would trigger immediate real-sector collapse. Accommodating the deficits of socialized enterprises becomes an institutional imperative to maintain economic reproduction.

\textbf{3.14} This reinterpretation establishes the central theoretical sequence of the chapter:
\[
\begin{gathered}
\text{financial fragility} \;\longrightarrow\; \text{peripheral external amplification} \\
\longrightarrow\; \text{central-bank balance-sheet exposure} \;\longrightarrow\; \text{state-dependent accommodation}.
\end{gathered}
\]
The macroeconomic transmission behavior of the peripheral central bank is state-dependent on its balance-sheet solvency conditions. When foreign-exchange backing is adequate, the central bank possesses operational policy space: domestic credit expansion mobilizes idle capacity, and external price shocks are buffered at the border. As balance-sheet conditions deteriorate past critical thresholds, this buffer dissolves: defensive monetary emission, lacking hard-currency backing, transmits directly into price acceleration and currency depreciation. Crucially, the empirical regimes identified in this chapter are not derived from Minsky's categorical taxonomy; they are estimated directly from the data through the Solvency Growth Gap ($\Delta s_{t-1}$) in the Threshold Vector Autoregression (TVAR) developed in Section~\ref{sec:historical_empirical_results}.
